\section{Introduction}

Long-duration scalp electroencephalography (EEG) monitoring can support the
identification of epileptic seizures, but practical systems must control memory,
latency, and energy in addition to classification performance. Reported
CHB-MIT results are not a single leaderboard: they differ in task definition,
channel selection, labels, sampling, and patient split. In particular,
window-level accuracy cannot substitute for continuous event sensitivity,
false-alarm rate, and detection delay \cite{ali2024,shoeb2010}.

Hardware-oriented work has demonstrated that compact seizure networks can be
implemented on embedded processors and custom accelerators. Bahr \emph{et al.}
deployed a 10,162-parameter CNN on a GAP8 RISC-V processor, while Li \emph{et
al.} co-designed a compact parallel CNN and a memristive accelerator
\cite{bahr2021,li2022}. These studies motivate a hardware contribution that
preserves a fixed model contract and reports measured implementation evidence,
rather than treating parameter count as a proxy for deployment cost.

This work targets an AMD Kria KV260 FPGA platform and starts from a frozen
software reference, \episepset{}. The reference uses raw, normalized EEG input
and depthwise-separable temporal and pointwise spatial convolutions. Its small
static graph permits a controlled investigation of integer arithmetic,
buffering, and Conv1D parallelization on programmable logic.

The contributions are as follows:
\begin{enumerate}
    \item a frozen 17-channel, two-second EEG inference contract with a
    5,013-parameter FP32 reference and an INT16 tensor package;
    \item an integer-only verification methodology that separates software
    emulation, C reference execution, RTL co-simulation, and board execution;
    \item a KV260 implementation study reporting post-route resources,
    frequency, batch-one latency, throughput, power, energy, and model-output
    agreement; and
    \item a transparent separation between deployment evidence and biomedical
    validation evidence, preventing unsupported clinical claims.
\end{enumerate}

Items 1 and the software portion of item 2 are complete. The hardware evidence
in items 2--3 is reported only after the planned KV260 experiments.
